\begin{homeworkProblem}[Seção 3-5]
  Seja $f:A\to\mathbb{R}$ uma função limitada e não negativa no bloco $A\subseteq\mathbb{R}^{m}$. A fim de que $f$ seja integravél, é necessário e suficiente que o conjunto
  \begin{align*}
    C(f)=\{(x,y)\in\mathbb{R}^{m+1}:x\in A,0\leq y\leq f(x)\}
  \end{align*}
  seja $J$-mensurável. No caso afirmtivo, tem-se
  \begin{align*}
    \int_{A}f(x)\,dx=Vol(C(f)).
  \end{align*}
  \begin{solution}
  \textbf{Necessário:}\\
    Suponha que $f$ é integrável e vamos ver que $C(f)$ é $J$-mensurável.\\
    Seja $r>0$, logo para cada $x\in A$, defina o cubo de centro $x$ e aresta $r$ como $C(x,r)$, logo temos que $A\subseteq \bigcup_{x\in A}C(x,r)$, mas pelo teorema de Lindelöf, temos que existe uma subcoleção enumerável de cubos taes que $A\subseteq \bigcup_{i=1}^{\infty}C_i$.\\
    Assim, definamos
    \begin{align*}
      C_i(f)=\{(x,y)\in\mathbb{R}^{m+1}:x\in A\cap C_i\text{ e }0\leq y\leq f(x)\}.
    \end{align*}
    Agora, note que como $f$ é limitada no bloco $A$, temos que existe $M>0$ tal que $f(x)\in [0,M]$ para todo $x\in A$, logo temos que $C_i(f)\subseteq C_i\times [0,M]$, pelo que temos que $C_i(f)$ é limitado. Portanto, para ver que $C_i(f)$ é J-mensurável, basta ver que $med(Fr(C_i(f)))=0$.\\
    Então, nos podemos ver o seguinte
    \begin{align*}
      Fr(C_i(f))=&\{(x,y)\in \mathbb{R}^{m+1}:x\in A\cap C_i\text{ e }y\in\{0,f(x)\}\text{ ou }x\in D_f\cap C_i\text{ e }y\in\mathbb{R}\}\\
      =&A\cap C_i\times \{0\}\cup \{(x,f(x)):x\in A\cap C_i\}\cup D_f\cap C_i\times \mathbb{R}\\
      :=& I+J+K.
    \end{align*}
    \begin{itemize}
      \item Note que $med(I)=0$, ja que como $\{0\}$ é um conjunto de medida nula no $\mathbb{R}$, então $A\cap C_i\times \{0\}$ é um conjunto de medida nula no $\mathbb{R}^{m+1}$.
      \item Note que $med(K)=0$, ja que como $D_f\cap C_i\subseteq D_f$ e $med(D_f)=0$ (pois $f$ é integrável), então $med(D_f\cap C_i)=0$.
      \item Vamos ver que $med(J)=0$, note que
        \begin{align*}
          J=\{(x,f(x)):x\in A\cap C_i\cap D_f\}\cup\{(x,f(x)):x\in A\cap C_i\setminus D_f\}:=J_1+J_2.
        \end{align*}
        Logo note que como $A\cap C_i\cap D_f\subseteq D_f$ e $med(D_f)=0$, então como temos que $J_1\subseteq D_f\times [0,M]$, podemos dizer que $med(J_1)=0$.\\
        Agora, vamos ver que $J_2$ tem medida nula.\\
        Note que como $x\in \tilde{J}=A\cap C_i\setminus D_f$, temos que $f$ é contínua no $\tilde{J}$, isto é que dado $\epsilon>0$ e dado $x\in \tilde{J}$, existe $\delta>0$ tal que se $y\in C(x,\delta)$ (cubo de centro $x$ e aresta $\delta$) então $f(y)\in C(f(x),\epsilon)$.\\
        Assim, observe que dado $\epsilon>0$, temos que para cada $x\in \tilde{J}$ existe $\delta_x>0$ tal que se $y\in C(x,\delta_x)\cap \tilde{J}$, então $f(y)\in C(f(x),\epsilon)$, logo note que $\tilde{J}\subseteq \bigcup_{x\in \tilde{J}}C(x,\delta_x)$, mas como $\tilde{J}\subseteq C_i$ e $\bigcup_{x\in \tilde{J}}C(x,\delta_x)$ é fechado, temos que $\bigcup_{x\in\tilde{J}}C(x,\delta_x)\cap C_i$ é fechado e limitado, logo compacto, pelo que sabemos que existe $k\in\mathbb{Z}^{+}$ tal que $\tilde{J}= \bigcup_{j=1}^{k}C_j$ com $C_j=C(x_j,\delta_{x_j})\cap \tilde{J}$, logo temos que
        \begin{align*}
          med(J_2)\leq \sum_{j=1}^{k}med(C_{j}\times C(f(x_j),\epsilon))
          \leq \sum_{j=1}^{k}med(C_j)\epsilon
          \leq \epsilon \sum_{j=1}^{k}med(C_j)
          \leq \epsilon\, med(\tilde{J}).
        \end{align*}
        Isto para todo $\epsilon>0$, logo podemos concluir que $med(J_2)=0$, depois $med(J)=0$. 
    \end{itemize}
    Assim, temos que $med(Fr(C_i(f)))=0$ e portanto que $C_i(f)$ é J-mesurável.\\
    Agora, note que $A=\bigcup_{i=1}^{\infty}C_i\cap A$, logo
    \begin{align*}
      C(f)=\bigcup_{i=1}^{\infty}C_i(f)=\bigcup_{i=1}^{\infty}\{(x,y)\in\mathbb{R}^{m+1}:A\cap C_i,\text{ e }0\leq y\leq f(x)\}.
    \end{align*}
    Portanto, como união numerável de conjuntos J-memsuravéis é J-mensuravél, temos que $C(f)$ é J-mensurável.\\
  \textbf{Suficiente:}\\
    Suponha que $C(f)$ é J-mensurável e vamos ver que $f$ é integrável.\\
    Note que $f$ é uma função limitada e não negativa no bloco $A\subseteq\mathbb{R}^{m}$, então existe um $M>0$ tal que $C(f)\subseteq A\times [0,M]$, logo definindo $\xi:A\times[0,M]\to\mathbb{R}$ dada por
    \begin{align*}
      \xi:A\times[0,M]&\to\mathbb{R}\\
                    (x,y)&\to\begin{cases}
                      1, &\text{ se } y\leq f(x) \text{,} \\
                      0, &\text{ se } y>f(x).
                    \end{cases}
    \end{align*}
    Temos que como $C(f)$ é J-mensurável, então $\xi$ é integrável e
    \begin{align*}
      \int_{A\times[0,M]}\xi(z)\,dz= Vol(C(f)).
    \end{align*}
    Note que como $\xi$ é integrável, temos que dado $\epsilon>0$ existe uma partiçao $P=P_1\times P_2$ com $P_1$ partição de $A$ e $P_2$ partição de $[0,M]$ tal que
    \begin{equation}\label{eq:S}
      \sum_{\substack{B_1\in P_1\\ B_2\in P_2}}M_{B_1\times B_2}Vol(B_1)Vol(B_2)-Vol(C(f))<\epsilon.
    \end{equation}
    Sem perdida de generalidad, suponja que $\displaystyle \max_{x\in B_1}f(x)\in P_2$ e $\displaystyle \min_{x\in B_1}f(x)\in P_2$ para cada $B_1\in P_1$, isto ja que se pegamos $\tilde{P}_2=P_2\cup\{\max_{x\in B_1}f(x),\min_{x\in B_1}f(x):B_1\in P_1\}$, então $P_2\subseteq \tilde{P}_2$, pelo que temos que $S(\xi, P_1\times P_2)\geq S(\xi,P_1\times\tilde{P}_2)$, logo temos que se $B_2\in \tilde{P}_2$, então a equação \eqref{eq:S} se cumpre.\\
    Agora, note que como $\max_{x\in B_1}f(x)\in P_2$, temos que ou $B_2\subseteq [0,\max_{x\in B_1}f(x)]$ ou $B_2\cap [0,\max_{x\in B_1}f(x)]=\emptyset$, i.é, que o $\min\{y\in B_2\}> \max_{x\in B_1}f(x)$, logo
    \begin{align*}
      M_{B_1\times B_2}=\max_{(x,y)\in B_1\times B_2}\xi(x,y)= 
      \begin{cases}
        1, &\text{ se } \displaystyle B_2\subseteq \left[0,\max_{x\in B_1}f(x)\right]\text{,} \\
        0, &\text{ se } \displaystyle B_2\cap \left[0,\min_{x\in B_1}f(x)\right]=\emptyset.
      \end{cases}
    \end{align*}
    Logo como $\displaystyle \max_{x\in B_1}f(x)\in P_2$, então temos que
    \begin{align*}
      \sum_{\substack{B_1\in P_1\\ B_2\in P_2}}M_{B_1\times B_2}Vol(B_1)Vol(B_2)=&\sum_{ B_1\in P_1}\left[\sum_{B_2\cap \left[0,\max_{x\in B_1}f(x)\right]\neq\emptyset}Vol(B_2)\right]Vol(B_1)\\
      =&\sum_{B_1\in P_1}\max_{x\in B_1}f(x)Vol(B_1)\\
      =&S(f,P_1).
    \end{align*}
    Agora, note que da mesma maneira temos que como $\min_{x\in B_1}f(x)\in P_2$, temos que ou $B_2\subseteq [0,\min_{x\in B_1}f(x)]$ ou $B_2\cap [0,\min_{x\in B_1}f(x)]=\emptyset$, i.é, que o $\min\{y\in B_2\}> \min_{x\in B_1}f(x)$, logo
    \begin{align*}
      m_{B_1\times B_2}=\min_{(x,y)\in B_1\times B_2}\xi(x,y)=\begin{cases}
        1, &\text{ se } B_2\subseteq\displaystyle \left[0,\min_{x\in B_1}f(x)\right]  \text{,} \\
        0, &\text{ se } B_2\cap \displaystyle\left[0,\min_{x\in B_1}f(x)\right]=\emptyset.
      \end{cases}
    \end{align*}
    Asimm,
    \begin{align*}
      \sum_{\substack{B_1\in P_1\\ B_2\in P_2}}m_{B_1\times B_2}Vol(B_1)Vol(B_2)=&\sum_{B_1\in P_1}\left[\sum_{B_2\subseteq \left[0,\min_{x\in B_1}f(x)\right]\neq\emptyset}Vol(B_2)\right]Vol(B_1)\\
      =&\sum_{B_1\in P_1}\min_{x\in B_1}f(x)Vol(B_1)\\
      =&s(f,P_1).      
    \end{align*}
    Logo como pela caracterização de funções integráveis temos que se $S(\xi,P)-s(\xi,P)<\epsilon$, então $S(f,P_1)-s(f,P_1)<\epsilon$, pelo que podemos afirmar que $f$ é uma função integrável.\\
    Além mais, pela equação \eqref{eq:S} temos que
    \begin{align*}
      S(f,P_1)-Vol(C(f))<\epsilon.
    \end{align*}
    Pelo que podemos concluir que
    \begin{align*}
      \int_{A}f(x)\,dx=Vol(C(f)).
    \end{align*}
  \end{solution}
\end{homeworkProblem}
